% Athena in a Nutshell
% !TEX root = main.tex

\chapter{Athena in a Nutshell}\label{chap:athena}

This chapter introduces Athena, the declarative proof language used over this thesis for formalizing CRDT specifications. Athena~\cite{athenaref} is a multi-paradigm language that combines automated reasoning with user-guided proof construction, making it well-suited for theorem proving about algebraic structures like semilattices.

We present the essential Athena constructs through practical examples, building from basic declarations to complete proofs. This prepares readers to understand the formal specifications without requiring prior Athena experience.

\section{Why Athena for CRDT Verification?}

CRDTs rely on algebraic properties (commutativity, associativity, idempotence) that must hold for all states and operations. Athena provides:

\begin{itemize}
  \item \textbf{Equational reasoning}: Natural expression properties through equality chains
  \item \textbf{Modularity}: Theory hierarchies that factor common proofs
  \item \textbf{Automation}: Built-in inference methods reduce proof load.
\end{itemize}

\section{Basic Constructs}

\subsection{Modules and Namespaces}

Modules encapsulate related definitions and theorems. Each module creates a namespace that enables a hierarchical organization of proofs.:

\begin{lstlisting}[language=Athena,style=athena,caption={Module definition with nested namespace}]
module GCounter {
  # Declarations visible within GCounter module
  declare merge: [N N] -> N
  
  module Theory {
    # Nested module for theorems
    conclude commutative := (forall ?n1:N ?n2:N . 
      (merge ?n1 ?n2) = (merge ?n2 ?n1))}
}

(!chain [(GCounter.merge x y) 
       = (GCounter.merge y x)  [GCounter.Theory.commutative]])
\end{lstlisting}

The \lstinline{open} directive imports module contents into the current namespace, reducing verbosity when working extensively with a module's statements.

\subsection{Sorts and Declarations}

Athena is multi-sorted: every term has a sort (type), and operations declare their input and output sorts. The \lstinline{declare} keyword introduces uninterpreted function symbols:

\begin{lstlisting}[language=Athena,style=athena,caption={Sort declarations and function signatures}]
# Define custom sorts
define-sort State := (Pair N N)  # Product type
datatype Value := (int N) | (str String)  # Sum type

# Declare operations with sort signatures
declare initial: [] -> State        # Nullary (constant)
declare update: [State Value] -> State  # Binary function
declare merge: [State State] -> State   # Binary operation

define [s s1 s2] := [?s:State ?s1:State ?s2:State]
\end{lstlisting}

Pattern variables (prefixed with \lstinline{?}) range over their declared sorts. The expression \lstinline{define [s ...] := [?s:State ...]} reduces repetition in axioms and theorems.

\subsection{Assertions: Stating Axioms}

The \lstinline{assert} directive introduces axioms/propositions assumed without proof. The \lstinline{assert*} variant allows multiple axioms in a list:

\begin{lstlisting}[language=Athena,style=athena,caption={Axiomatic specification of merge}]
# Single assertion
assert merge-commutative := 
  (forall ?s1:State ?s2:State . (merge s1 s2) = (merge s2 s1))
# Multiple related axioms
assert* merge-axioms := [
  (forall ?s1 ?s2 . (merge s1 s2) = (merge s2 s1))  # Commutativity
  (forall ?s . (merge s s) = s)                      # Idempotence]
\end{lstlisting}

Axioms form the foundation of a specification.

\subsection{Proofs: The \texttt{conclude} Construct}

While \lstinline{assert} states axioms, \lstinline{conclude} proves theorems. A proof provides a \emph{deduction} that justifies the conclusion:

\begin{lstlisting}[language=Athena,style=athena,caption={Basic proof structure}]
conclude merge-symmetric := 
  (forall ?s1:State ?s2:State . 
    (merge s1 s2) = (merge s2 s1) ==> (merge s2 s1) = (merge s1 s2))
  pick-any s1:State s2:State
    assume hyp := ((merge s1 s2) = (merge s2 s1))
      (!sym hyp)  # Symmetry of equality
\end{lstlisting}

The proof structure follows the logical form:
\begin{enumerate}
  \item \lstinline{pick-any} introduces universally quantified variables
  \item \lstinline{assume} binds the hypothesis of an implication
  \item Inference methods (prefixed with \lstinline{!}) construct the deduction
\end{enumerate}

\section{Proof Methods}

Athena provides primitive inference rules as well as high-level proof methods. Methods automate common reasoning patterns.

\subsection{Equational Reasoning with \texttt{!chain}}

The \lstinline{!chain} method builds a chain of equalities, where each step is justified by an axiom or theorem:\\

\begin{lstlisting}[language=Athena,style=athena,caption={Equational chain proof}]
conclude merge-associative-demo := 
  (forall ?s1:State ?s2:State . 
    (merge (merge s1 s2) s2) = (merge s1 s2))
  pick-any s1:State s2:State
    (!chain [(merge (merge s1 s2) s2)
           = (merge s1 (merge s2 s2))  [merge-associative]
           = (merge s1 s2)             [merge-idempotent]])
\end{lstlisting}

Each line in the chain states an equality, and the term in brackets names the justification. This style emulates standard mathematical practice and makes proofs highly readable.

\subsection{Implication Chain with \texttt{!chain->}}

For proving implications, \lstinline{!chain->} sequences modus ponens steps:

\begin{lstlisting}[language=Athena,style=athena,caption={Implication chain}]
# Given: (P ==> Q) and (Q ==> R) axioms
conclude transitive-implication := (P ==> R)
  (!chain-> [P ==> Q  [axiom-1]
               ==> R  [axiom-2]])
\end{lstlisting}

This is particularly useful for monotonicity proofs, where we show that one property implies another through a sequence of intermediate steps.

\subsection{Universal Instantiation with \texttt{!uspec}}

When an axiom or theorem is universally quantified, \lstinline{!uspec} specializes it to concrete terms:

\begin{lstlisting}[language=Athena,style=athena,caption={Universal specialization}]
# Given axiom: (forall ?s1 ?s2 . (merge s1 s2) = (merge s2 s1))
assert merge-commutative := (forall ?s1 ?s2 . ...)

# Specialize to specific states
define instance := (!uspec merge-commutative s1 s2)
# Now instance is: (merge s1 s2) = (merge s2 s1)
\end{lstlisting}

Athena offers both \lstinline{!chain} and \lstinline{!uspec} as formal methods. The following are equivalent:

\begin{lstlisting}[language=Athena,style=athena,caption={Automatic vs manual instantiation}]
# Manual
define step1 := (!uspec merge-commutative s1 s2)
define result := (!chain [(merge s1 s2) = (merge s2 s1) [step1]])

# Automatic
define result := (!chain [(merge s1 s2) = (merge s2 s1) [merge-commutative]])
\end{lstlisting}

\subsection{Case Analysis with \texttt{!cases}}

For datatype definitions, \lstinline{!cases} performs structural induction or case splits:

\begin{lstlisting}[language=Athena,style=athena,caption={Proof by cases over a datatype}]
datatype Update := (increment) | (decrement)

conclude update-total := 
  (forall ?u:Update . (u = increment) | (u = decrement))
  pick-any u:Update
    (!cases u {
      (increment) => (!left (!reflex increment))
      (decrement) => (!right (!reflex decrement)) })
\end{lstlisting}

\section{Theories and Adapters}

Athena's \emph{theory mechanism} enables modular proofs. A theory packages axioms and derived theorems, which can then be instantiated for multiple concrete types.

\subsection{Defining a Theory}

A theory is created with \lstinline{make-theory}, specifying axioms and theorems:

\begin{lstlisting}[language=Athena,style=athena,caption={Abstract semilattice theory}]
module SemilatticeTheory {
  declare join: (T) [T T] -> T
  assert commutative := (forall ?x ?y . (x join y) = (y join x))
  assert associative := (forall ?x ?y ?z . 
    ((x join y) join z) = (x join (y join z)))
  assert idempotent := (forall ?x . (x join x) = x)
  
  # Package into a theory
  define theory := (make-theory [] [commutative associative idempotent])
}
\end{lstlisting}

\subsection{Instantiating via Adapters}

To apply the theory to a concrete CRDT, we create an \emph{adapter} that maps the abstract operations to concrete ones:

\begin{lstlisting}[language=Athena,style=athena,caption={CRDT adapter inheriting semilattice properties}]
module GCounterAdapter {
  # Map abstract 'join' to concrete 'gc-merge'
  define merge := gc-merge
  
  # Prove adapter satisfies theory axioms
  conclude commutative := 
    (forall ?s1:N ?s2:N . (gc-merge s1 s2) = (gc-merge s2 s1))
    pick-any s1:N s2:N
      (!chain [(gc-merge s1 s2)
             = (N.max s1 s2)          [gc-merge-definition]
             = (N.max s2 s1)          [N.Max.commutative]
             = (gc-merge s2 s1)       [gc-merge-definition]])
  
  # Similar for associative, idempotent...
  # Now all theorems from SemilatticeTheory hold for gc-merge
}
\end{lstlisting}

This pattern eliminates redundant proofs: once we've proven that \lstinline{N.max} is a semilattice, any CRDT using max-based merge inherits those properties.

\section{Reading Athena Code}

When encountering Athena code in later chapters, readers should:

\begin{enumerate}
  \item \textbf{Identify the module structure}: What operations are declared? What sorts are involved?
  \item \textbf{Locate the axioms}: Which properties are asserted vs. proven?
  \item \textbf{Follow the proof chains}: Equational steps mirror standard algebraic reasoning
  \item \textbf{Check the justifications}: Each step cites an axiom, theorem, or inference rule
\end{enumerate}

The formal specifications serve two purposes:
\begin{itemize}
  \item \textbf{Documentation}: Precise semantics for CRDT operations
  \item \textbf{Verification}: Machine-checked proofs of convergence properties
\end{itemize}

\section{Tools and Workflow}

Athena proofs are developed interactively. A typical workflow:

\begin{enumerate}
  \item Write the specification (sorts, declarations, axioms)
  \item State the theorem to prove with \lstinline{conclude}
  \item Build the proof incrementally, checking each step
  \item Refactor common patterns into reusable methods
\end{enumerate}

The Athena interpreter provides immediate feedback on type errors, failed deductions, and missing justifications, enabling rapid iteration.

\bigskip

With these fundamentals in hand, readers are prepared to understand the formal CRDT specifications in Chapter~\ref{chap:theory} and the framework implementations in Chapter~\ref{chap:framework}. The subsequent chapters assume familiarity with the definitions presented here, focusing on the algebraic structures and convergence properties rather than Athena syntax.

\endinput
